% Verification protocol and threats to validity.
\section{Verification protocol and threats to validity}\label{app:verification}

\subsection{Protocol}

Every claim in the paper rests on one or more of the following checks.

\begin{enumerate}
\item \textbf{Solution files.} Every solution file in the data package is
written in BOBILib's format and passes \texttt{bobi verify}
(Section~\ref{app:notation}): all rows, bounds and integrality requirements of
the HPR within $10^{-6}$; the stated leader objective; and the follower's
problem with the leader fixed, re-solved by the exact knapsack oracle (one
follower row) or by a HiGHS MILP built from the model (otherwise), accepted
by BOBILib's relative test~\eqref{eq:bobilib-tol}. This covers the 583
imported KIP files, the 16 subset-sum files, the 690 OR files (270 with one
follower knapsack, 420 with several), the 64 general files, the 29 files with
coupling rows, the 60 clique files and the 36 network files. The points with
coupling rows also pass \texttt{bobi certify} (proven dual bounds, relative
$10^{-5}$, optimistic tie-break).

\item \textbf{Optimality proofs.} Each ``optimal'' status rests on a lower
bound of one of these kinds: a published proof (cited, with the published
value matched by a verified solution file); the exact subset-sum search
(Section~\ref{app:kip}), whose pruning rules are proved in Lean; a theorem
with a checkable certificate (Mantel, Brouwer and partitions for clique
interdiction, Section~\ref{app:clique}; the HPR bound with an exact
model comparison, Section~\ref{app:general}); or exhaustive enumeration with
checked witnesses (Section~\ref{app:network}). ``No bilevel feasible point''
rests on Theorem~\ref{thm:ray} with an exact rational check (38 instances)
or a floating-point check (2 instances).

\item \textbf{Independent provers for the subset-sum proofs.} Two Python
provers that share no code with the Rust search, one written by the
designer and one by a reviewer who read neither, parse the instance files
themselves and re-prove all 16 lower bounds with only the direct rule and the
per-node rule (Theorem~\ref{thm:node-rule}), in integer arithmetic; both
compare themselves with brute force on random small instances. Two other
designs of the search, built separately, prove
\texttt{interKP-200-100-6-3}, and one of them all four subset-sum instances
with budget index~3.
The reviewer compared every search binary with exhaustive enumeration on
2,990 random small instances (no mismatch).

\item \textbf{Tests of the code against its specification.} Unit tests
compare the subset-sum search with brute force over every deletion, the
threshold tables with the bitset recursion, the node tables with the online
game defined directly and with enumeration of every deletion, and the
table code with a direct transcription of the Lean recursion
$\mathrm{thr}$ (also with costs not ascending, as after restricted copies).
The network DP is compared with brute force over all follower points on
random forests, and the whole enumeration with a brute-force bilevel
optimum.

\item \textbf{Independent reviews.} The clique, network and coupled-row
results were each re-derived by a separate reviewer with its own instance
reader and different solvers or code: two exact maximum-clique searches,
HiGHS and a branch and bound written in C for the part-size cases of
Proposition~\ref{prop:clique-sizes}; a C implementation of the tree DP and a
brute force based on Lemma~\ref{lem:fire-closure}, compared with ours at all
9.3 million attacks; and a second reading of the coupled files with HiGHS's
own MPS reader. The reviewers found only wording errors.

\item \textbf{Exact arithmetic where it is cheap.} The clique inequalities
(Mantel, Brouwer's threshold, the size vectors of
Proposition~\ref{prop:clique-sizes}), the recession directions and the
comparison of each HPR with its MIPLIB model are checked with integers or
exact rationals.

\item \textbf{Machine-checked lemmas.} The Lean~4 project (with Mathlib)
proves the lemmas listed in Section~\ref{app:lean-map}. It contains no
\texttt{sorry}, and \texttt{\#print axioms} shows only Lean's standard axioms
(\texttt{propext}, \texttt{Classical.choice}, \texttt{Quot.sound}) for each
main theorem.
\end{enumerate}

\subsection{Threats to validity}

\begin{enumerate}
\item \textbf{Trusted solvers.} Several verdicts come from MILP or CP solvers:
follower optimality for followers with several rows and for general
instances (HiGHS), the 170 infeasibility verdicts behind $w_3(G)>k$ (each
reproduced by five methods, two of them solver-free), the follower optimum
at the minimising attack in the network proof (HiGHS), and the dual bounds
in the OR bracket (Theorem~\ref{thm:mkip-bracket}). A solver bug that all
methods share would go unnoticed.

\item \textbf{Floating point and tolerances.} Points are feasible within
$10^{-6}$ and follower-optimal within BOBILib's relative $10^{-4}$, not
exactly. Nine coupled points use that slack, and two coupled files are
feasible only within it. The timtab1 point lies $2.2\times10^{-7}$ below its
bound. Two of the 40 ray certificates are checked only in floating point.
The merged-row online bound (Theorem~\ref{thm:mkip-omega}) and the LP
bracket are evaluated in double precision, so their rounding steps are in
the safe direction only up to floating-point error. The window rule of the
subset-sum search uses a floating-point fractional bound; the 16 proofs are
independently confirmed by the integer-only provers above, which do not use
that rule.

\item \textbf{Scope of the Lean proofs.} Lean proves lemmas about abstract
models, not the code. That the Rust loops compute the recursions, the range
shortcuts and early exits of the table code, the bitset and cache code, and
the absence of integer overflow are argued in doc comments and tested, not
formalised. That the instance files have the stated parameters ($n$, $m$, $k$,
budgets) is checked by scripts, not by Lean. Section~\ref{app:lean-map} lists
what is not formalised.

\item \textbf{Results taken from the literature.} Brouwer's theorem was read
in the scanned 1981 report and agrees with later restatements, but it is
neither formalised nor reproved here. Published optima (Weninger and
Fukasawa for KIP, Fischetti et al.\ for imkp, Lozano, Bergman and Cire for
OR, MIPLIB optima for the HPR bounds) and the published lower bounds of the
OR instances with several follower knapsacks are used as reported; for
them only the points are verified here.

\item \textbf{Instance reading.} Every argument depends on reading the MPS
and AUX files correctly. Each family's model is checked row by row against
its expected pattern, and the key results were re-derived with independent
readers (two subset-sum provers, the exact ray and HPR scripts, the
reviewers' parsers).

\item \textbf{Running times.} All times were measured on a shared machine
with four cores that was often oversubscribed; they are indicative only and
do not affect any proof.
\end{enumerate}
